\title{Controlling Text-to-Image Diffusion by Orthogonal Finetuning}

\begin{abstract}
Text-to-image diffusion models at scale have demonstrated remarkable proficiency in synthesizing realistic images from textual descriptions. However, steering these sophisticated models to address a variety of downstream objectives remains a significant unresolved issue. To confront this, we propose a systematic finetuning strategy termed Orthogonal Finetuning (OFT), designed for tailoring text-to-image diffusion models to specific tasks. Distinct from previous approaches, OFT offers theoretical guarantees for the preservation of hyperspherical energy, a metric representing inter-neuronal relationships on the unit hypersphere. Our investigations reveal that maintaining this property is vital for sustaining the semantic generative capacity of these models. To enhance stability throughout the finetuning process, we introduce Constrained Orthogonal Finetuning (COFT), which applies an extra radial constraint on the hypersphere. We focus on two critical downstream finetuning tasks: subject-driven generation, which entails generating images centered on a particular subject given a handful of reference images and a textual description; and controllable generation, in which the model is conditioned to accommodate supplementary control inputs. Empirical results indicate that our OFT framework surpasses existing finetuning methodologies with respect to both output quality and convergence efficiency.
\end{abstract}

\section{Introduction}

Contemporary text-to-image diffusion frameworks~\cite{saharia2022photorealistic,ramesh2022hierarchical,rombach2022high} deliver state-of-the-art outcomes in synthesizing high-quality images from text instructions. Nonetheless, reliance solely on textual input often leads to imprecise or vague control over generated visual content. In response, our study is concerned with two central text-to-image generation challenges:

\begin{itemize}[leftmargin=*,nosep]

    \item \textbf{Subject-driven generation}~\cite{ruiz2023dreambooth}: Provided with several reference images of a particular subject, the objective is to render novel images depicting the same subject in various scenarios, steered by accompanying text prompts.
    \item \textbf{Controllable generation}~\cite{zhang2023adding,mou2023t2i}: When supplied with an additional conditioning cue (such as canny edges or segmentation maps), the task is to create images that conform to both the provided control input and the textual directive.
\end{itemize}

The core challenge in both tasks lies in how to fine-tune text-to-image diffusion models in a manner that maintains the strong generative capabilities established during pretraining. An effective finetuning approach should satisfy two main criteria: (1) \emph{training efficiency}, characterized by reduced numbers of adjustable parameters and minimized training epochs, and (2) \emph{generalizability preservation}, meaning the model retains its ability to generate high-quality and diverse outputs. Typically, finetuning methods address these goals by either making incremental updates to neuron weights using small learning rates (\emph{e.g.}, \cite{ruiz2023dreambooth}) or by incorporating lightweight modules through re-parameterized neuron weights (\emph{e.g.}, \cite{hulora2022,zhang2023adding}). However, despite their straightforward nature, such strategies do not inherently guarantee that the pretrained generative performance will be preserved. Moreover, there is a notable absence of theoretical frameworks for systematically developing finetuning algorithms or selecting appropriate hyperparameters such as training duration. One significant challenge is the lack of robust metrics for evaluating the retention of generative abilities after pretraining. Current approaches generally assume that a reduced Euclidean norm between finetuned and pretrained models corresponds to better preservation of the original capabilities. As a result, very conservative learning rates are commonly used during finetuning. Nevertheless, this assumption is limited; Euclidean distance alone does not adequately reflect the extent of semantic retention. Consequently, we contend that more nuanced, structurally informed metrics are essential for more accurately assessing the changes introduced by finetuning, which would, in turn, lead to improved preservation of generative performance and increased stability during the finetuning process.

Building on empirical findings that hyperspherical similarity effectively represents semantic relationships~\cite{liu2017deep,liu2018decoupled,chen2020angular}, our approach leverages hyperspherical energy~\cite{liu2018learning} to model the pairwise associations among neurons. This quantity, computed as the aggregate hyperspherical similarity (such as cosine similarity) across all pairs of neurons within a layer, reflects the uniformity of neuron distribution on the unit hypersphere~\cite{liu2021learning}. We posit that an optimal finetuned model should display minimal alteration in hyperspherical energy relative to its pretrained counterpart. One straightforward strategy would be to implement a regularization term to maintain consistent hyperspherical energy throughout finetuning; however, such an approach does not guarantee effective minimization of the energy discrepancy. To address this limitation, we utilize the invariance characteristic of hyperspherical energy—specifically, that applying a uniform orthogonal transformation to all neurons preserves their pairwise hyperspherical similarities. This observation underpins our method, Orthogonal Finetuning~(OFT), which adapts large text-to-image diffusion models for downstream tasks while ensuring the hyperspherical energy remains unchanged. The core mechanism involves learning an orthogonal transformation, shared across the layer, that maintains the angles between neuron vectors. Alternatively, OFT can be interpreted as modifying the standard coordinate basis for the neurons within a layer. Additionally, since closer Euclidean proximity between finetuned and pretrained models tends to preserve pretraining effectiveness, we introduce a variant named Constrained Orthogonal Finetuning~(COFT). This method enforces that the updated model resides on a hypersphere of constant radius, centered at the positions of the pretrained neurons.

The rationale behind the effectiveness of orthogonal transformation for neuron finetuning is partly inspired by the principles of 2D Fourier transforms, which allow images to be represented in terms of magnitude and phase spectra. Among these, the phase spectrum—capturing the angular relationship between the input and basis—retains the core semantic information. Notably, combining the phase spectrum from an image with a random magnitude spectrum enables a reconstruction that largely preserves the original image’s semantic content. This observation points to the pivotal role of neuron directions in modulating the semantics of generated images, aligning with the objectives of subject-driven and controllable image generation. Nevertheless, allowing neuron directions to change without restriction could severely compromise the generative performance established during pretraining. To mitigate this, we introduce a constraint to preserve the pairwise angles between neurons, grounded in the hypothesis that these angular relationships encode much of the neural network's embedded knowledge. Building on this foundation, we advocate learning a layer-shared orthogonal transformation for neurons at each layer, thereby ensuring that the hyperspherical energy remains constant.

Our methodology also draws on orthogonal over-parameterized training~\cite{liu2021orthogonal}, wherein classification neural networks are trained from scratch by orthogonally transforming randomly initialized networks. Since such random initializations are expected to yield networks with low hyperspherical energy, the central objective of~\cite{liu2021orthogonal} is to maintain this low energy throughout training—an approach that correlates with improved generalization in classification tasks~\cite{liu2018learning,lin2020regularizing}. Their results demonstrate that orthogonal transformation affords sufficient expressiveness for training neural networks to generalize well in classification contexts. However, our focus differs as we address the finetuning of text-to-image diffusion models to enhance both control and generative performance on downstream tasks. We emphasize two primary distinctions between OFT and the approach in~\cite{liu2021orthogonal}. First, while~\cite{liu2021orthogonal} targets the minimization of hyperspherical energy, OFT is designed to maintain the hyperspherical energy at the level established by the pretrained model, thus safeguarding the original internal structure during finetuning. In contrast, minimizing hyperspherical energy in diffusion models may degrade the fidelity of their semantic structures. Second, OFT explicitly constrains deviation from the pretrained model, resulting in a constrained variant, whereas~\cite{liu2021orthogonal} does not implement such constraints. Effective finetuning thus relies on balancing adaptability with stability, a trade-off our OFT framework is specifically designed to optimize. Our contributions are summarized as follows:

\begin{itemize}[leftmargin=*,nosep]

    \item We introduce Orthogonal Finetuning, a new approach to finetuning that enhances the controllability of text-to-image diffusion models. Additionally, we propose a constrained variant of this method that restricts angular deviations from the original pretrained model, which serves to further stabilize the finetuning process.
    \item In contrast with prevailing finetuning techniques, OFT guarantees the preservation of hyperspherical energy during finetuning—a property we empirically identify as critical for maintaining the semantic generative capacities of the pretrained model.
    \item We demonstrate the utility of OFT in both subject-driven and controllable generation settings. Through an extensive suite of experiments, we show that OFT consistently yields substantial gains over previous methods in terms of sample quality, speed of convergence, and robustness throughout finetuning. Furthermore, OFT exhibits enhanced sample efficiency, as it achieves high performance with considerably fewer training images and epochs.
    \item For tasks involving controllable generation, we present a novel control consistency metric to quantitatively assess how well controllability is maintained. The key procedure involves inferring the control signal from the produced image and juxtaposing it with the original control input, thereby affirming OFT’s superior controllability.
\end{itemize}

\section{Related Work}

\textbf{Text-to-image diffusion models}. Recent years have witnessed rapid advances in text-to-image synthesis~\cite{saharia2022photorealistic,ramesh2022hierarchical,rombach2022high,gu2022vector,nichol2022glide}, primarily driven by the evolution of diffusion-based generative models~\cite{ho2020denoising,song2021score,song2021denoising,dhariwal2021diffusion} and progress in learning joint visual and linguistic representations~\cite{su2020vlbert,lu2019vilbert,radford2021learning,wang2021simvlm,li2022blip,li2023blip,alayrac2022flamingo,schuhmann2022laion}. For instance, GLIDE~\cite{nichol2022glide} and Imagen~\cite{saharia2022photorealistic} operate with diffusion in the pixel space: GLIDE co-trains its text encoder together with the diffusion model using text-image pairs, whereas Imagen leverages a fixed, pretrained text encoder. On the other hand, diffusion in latent spaces is the focus of models like Stable Diffusion~\cite{rombach2022high} and DALL-E2~\cite{ramesh2022hierarchical}; Stable Diffusion makes use of VQ-GAN~\cite{esser2021taming} to form a latent visual codebook, while DALL-E2 employs CLIP~\cite{radford2021learning} to create a unified embedding space for images and text. Besides diffusion models, other paradigms including generative adversarial networks~\cite{reed2016generative,zhang2017stackgan,xu2018attngan,li2019controllable} and autoregressive models~\cite{ramesh2021zero,ding2021cogview,wu2022nuwa,yuscaling} have been explored for this task. As a model-agnostic method, OFT can be integrated with any text-to-image diffusion model.

\textbf{Subject-driven generation}. Approaches such as \cite{avrahami2022blended,nichol2022glide} utilize user-specified masks as supplementary conditions to inhibit alterations to the primary subject. Inversion-based methods~\cite{choi2021ilvr,dhariwal2021diffusion,ramesh2022hierarchical,gal2022image} enable modifications to the surrounding context while maintaining the integrity of the main subject. Alternatively, \cite{hertz2022prompt} achieve both global and localized edits without the need for input masks. However, these strategies are generally insufficient for retaining detailed, identity-specific information. Pivotal Tuning~\cite{roich2022pivotal} mitigates this by adjusting a generator near an inverted latent code while applying identity-preserving regularization. Similarly, \cite{nitzan2022mystyle} constructs a person-specific generative face prior from multiple images. The method in \cite{casanova2021instance} offers instance-level diversity, but sometimes at the cost of losing fine details unique to that instance. DreamBooth~\cite{ruiz2023dreambooth}, on the other hand, uses a specialized token and a small set of subject images to finetune the diffusion model via reconstruction and class-specific prior preservation losses. While OFT also builds upon the DreamBooth pipeline, it differentiates

\textbf{Controllable generation}. Image-to-image translation is commonly framed as a controllable generation problem, with earlier approaches predominantly utilizing conditional generative adversarial networks~\cite{isola2017image,zhu2017toward,wang2018high,park2019semantic,choi2018stargan}. In addition to these, diffusion models have recently found application in image-to-image translation tasks~\cite{saharia2022palette,voynov2022sketch,wang2022pretraining}. The introduction of ControlNet~\cite{zhang2023adding} marks an advance in this area by adapting and finetuning a pretrained diffusion model to respond to additional control signals, resulting in notable improvements in controllable generation. Similarly, T2I-Adapter~\cite{mou2023t2i} represents another parallel development, which enhances the controllability of generated images by finetuning a pretrained diffusion model. In line with the methodological setups of \cite{zhang2023adding,mou2023t2i}, we utilize OFT to finetune pretrained diffusion models, achieving consistently enhanced controllability even with reduced amounts of training data and fewer finetuned parameters. Crucially, OFT introduces no supplementary computational costs during inference.

\textbf{Model finetuning}. The practice of finetuning large pretrained models for specific downstream tasks is now widespread~\cite{devlin2018bert,brown2020language,he2022masked}. Adaptation strategies, such as those explored in (\emph{e.g.}, \cite{hulora2022,houlsby2019parameter,pfeiffer2020adapterfusion}), are particularly well studied in the natural language processing community. Among them, LoRA~\cite{hulora2022} is highly related to OFT, leveraging a low-rank approach to updating model weights additively during finetuning. In contrast, OFT employs a layer-wise shared orthogonal transformation to multiplicatively adjust neuron weights, thereby analytically preserving the angles between neuron vectors within the same layer, which enhances training stability.

\section{Orthogonal Finetuning}

\subsection{Why Does Orthogonal Transformation Make Sense?}

We begin our exploration by explaining the rationale behind using orthogonal transformations in the finetuning of text-to-image diffusion models. This discussion is organized around two key questions: (1) the motivation for finetuning neuron angles (\emph{i.e.}, directional adjustments), and (2) the justification for employing orthogonal transformations specifically for angle finetuning.

To address the first question, we build on empirical findings from \cite{liu2018decoupled,chen2020angular}, which indicate that differences in angular features are effective indicators of semantic divergence. SphereNet~\cite{liu2017deep} demonstrates that constraining all neural activations to lie on the surface of a unit hypersphere preserves the model's capacity and can enhance generalization, suggesting that neuron directions can encapsulate essential information from the data. To further highlight the significance of angular information, we designed a simple experiment illustrated in Figure~\ref{fig:toy_ae}. We trained a conventional convolutional autoencoder on flower images. During training, the standard inner product was employed to compute the feature map, where $z$ is the convolution kernel $\bm{w}$'s output on the input $\bm{x}$ within the sliding window. For evaluation, we examined three alternative strategies for constructing the feature map: (a) the inner product used in training, (b) relying solely on magnitude, and (c) extracting only angular components. Results in Figure~\ref{fig:toy_ae} reveal that the angular representation of neurons alone is nearly sufficient to reconstruct the original images, whereas magnitude carries minimal useful information. It is important to note that cosine-based activation (c) was not involved during training, which exclusively utilized the inner product. These outcomes suggest that neuron directions predominantly encode the semantic content of input images, indicating that altering neuron directions is likely a more effective means for semantic manipulation.

Concerning the second question, one of the most straightforward techniques to adjust neuron directions is to apply a collective rotation or reflection across all neurons within the same layer, leading to an orthogonal transformation. While more complex angular transformations—such as individually rotating neurons by different amounts—could offer greater flexibility, we find that orthogonal transformation provides an optimal balance between adaptability and regularization. Additionally, \cite{liu2021orthogonal} substantiates that orthogonal transformations are sufficiently expressive for training neural networks. To substantiate our perspective, we conducted a regularization experiment within the controllable image synthesis setup of ControlNet~\cite{zhang2023adding}, with corresponding outcomes presented in Figure~\ref{fig:ortho}. The findings show that our standard OFT method demonstrates stable performance and achieves precise control upon completion of training (at epoch 20). By contrast, omitting the orthogonality constraint results in the inability to produce realistic images or effective control. This experiment confirms that the orthogonality constraint is integral to OFT’s effectiveness.

\subsection{General Framework}

Finetuning in its standard form generally employs gradient descent with a reduced learning rate to adjust the weights of a model, or in some cases, specific layers. The low learning rate implicitly limits how much the updated parameters $\bm{M}$ deviate from their original pretrained values $\bm{M}^0$, such that conventional finetuning can be interpreted as training under an implicit regularization term $\thickmuskip=2mu \medmuskip=2mu \| \bm{M} - \bm{M}^0\|$. While this approach is intuitively reasonable, it might still offer too much flexibility—particularly for large-scale models. To address this, LoRA imposes a further restriction by constraining the update to a low-rank structure, namely $\thickmuskip=2mu \medmuskip=2mu \text{rank}(\bm{M} - \bm{M}^0) = r'$, where $r'$ is chosen to be small. In contrast, OFT enforces similarity at the neuron-pair level, requiring $\thickmuskip=2mu \medmuskip=2mu \| \text{HE}(\bm{M})-\text{HE}(\bm{M}^0) \| = 0$, where $\text{HE}(\cdot)$ represents the hyperspherical energy associated with a given weight matrix.

To illustrate this, consider a fully connected layer defined as $\thickmuskip=2mu \medmuskip=2mu \bm{W} = \{\bm{w}_1, \cdots, \bm{w}_n\} \in \mathbb{R}^{d\times n}$, where each $\bm{w}_i\in\mathbb{R}^d$ signifies the $i$-th neuron and $\bm{W}^0$ holds the corresponding pretrained weights. The output vector for this layer, $\thickmuskip=2mu \medmuskip=2mu \bm{z} \in \mathbb{R}^n$, is calculated as $\thickmuskip=2mu \medmuskip=2mu \bm{z} = \bm{W}^\top \bm{x}$, with input $\thickmuskip=2mu \medmuskip=2mu \bm{x} \in \mathbb{R}^d$. The OFT method seeks to minimize discrepancies in hyperspherical energy between the fine-tuned and the original pretrained model, formalized by:

\begin{equation}\label{eq:oft_principle}
\footnotesize
\min_{\bm{W}} \left\| \text{HE}(\bm{W})-\text{HE}(\bm{W}^0)\right\|~~\Leftrightarrow~~\min_{\bm{W}} \bigg{\|} \sum_{i\neq j} \|\hat{\bm{w}}_i-\hat{\bm{w}}_j\|^{-1}-\sum_{i\neq j} \|\hat{\bm{w}}^0_i-\hat{\bm{w}}^0_j\|^{-1}\bigg{\|}
\end{equation}

Here, $\thickmuskip=2mu \medmuskip=2mu \hat{\bm{w}}_i := \bm{w}_i / \| \bm{w}_i \|$ denotes the $i$-th neuron after normalization, and the hyperspherical energy for $\bm{W}$ is given by $\thickmuskip=2mu \medmuskip=2mu \text{HE}(\bm{W}) := \sum_{i\neq j} \|\hat{\bm{w}}_i-\hat{\bm{w}}_j\|^{-1}$. It is straightforward to see that the minimum of Eq.~\eqref{eq:oft_principle} is zero; this occurs when $\bm{W}$ can be written as an orthogonal transformation of $\bm{W}^0$, i.e., $\thickmuskip=2mu \medmuskip=2mu \bm{W} = \bm{R} \bm{W}^0$ for some orthogonal matrix $\thickmuskip=2mu \medmuskip=2mu \bm{R} \in \mathbb{R}^{d\times d}$ (where $\bm{R}$ is a rotation if its determinant is $1$, or a reflection if $-1$). This principle underlies OFT: finetuning is carried

\begin{equation}\label{eq:oft_general}
    \footnotesize
    \bm{z}=\bm{W}^\top\bm{x}=(\bm{R}\cdot\bm{W}^0)^\top\bm{x},~~~\text{s.t.}~\bm{R}^\top\bm{R}=\bm{R}\bm{R}^\top=\bm{I}
\end{equation}

Here, $\bm{W}$ stands for the weight matrix that has been fine-tuned using OFT, while $\bm{I}$ represents the identity matrix. An overview of OFT is depicted in Figure~\ref{fig:block}. Analogous to the zero-initialization approach employed in LoRA, it is essential that OFT fine-tunes the model beginning precisely from $\bm{W}^0$. This requirement is met by initializing the orthogonal matrix $\bm{R}$ as an identity matrix, thus ensuring that the process starts with the pretrained weights. To maintain the orthogonality of $\bm{R}$ throughout training, differential orthogonalization techniques, such as those discussed in \cite{liu2021orthogonal,lezcano2019cheap}, may be leveraged. Later, we examine approaches to enforce this orthogonality constraint efficiently.

\subsection{Efficient Orthogonal Parameterization}\label{sect:eff_ortho}

While classic orthogonalization techniques like the Gram-Schmidt process are theoretically differentiable, their computational demands render them impractical for large-scale models~\cite{liu2021orthogonal}. To enhance efficiency, we instead employ the Cayley parameterization to construct the orthogonal matrix. Concretely, this approach expresses the orthogonal matrix as $\thickmuskip=2mu \medmuskip=2mu \bm{R}=(\bm{I}+\bm{Q})(I-\bm{Q})^{-1}$, where $\bm{Q}$ is skew-symmetric, i.e., $\thickmuskip=2mu \medmuskip=2mu \bm{Q}=-\bm{Q}^\top$. This parameterization has a minor caveat—only orthogonal matrices with determinant $1$, corresponding to the special orthogonal group, can be realized using the Cayley transform. Nonetheless, empirical results show that this restriction does not negatively impact performance. However, even with the Cayley parameterization, modeling a full orthogonal matrix $\bm{R}$ can become parameter-intensive when $d$ is large. To improve parameter efficiency, we advocate representing $\bm{R}$ as a block-diagonal matrix composed of $r$ blocks, resulting in the following structure:

\begin{equation}
    \footnotesize
    \bm{R}=\text{diag}(\bm{R}_1,\bm{R}_2,\cdots,\bm{R}_r)=\begin{bmatrix}
    \bm{R}_1\in \text{O}(\frac{d}{r}) &  &\\
    &\ddots & \\
    & & \bm{R}_r\in \text{O}(\frac{d}{r})
    \end{bmatrix}\in \text{O}(d)
\end{equation}

Here, $\text{O}(d)$ refers to the group of orthogonal matrices in $d$ dimensions, where $\thickmuskip=2mu \medmuskip=2mu \bm{R}\in\mathbb{R}^{d\times d}$ and, for all $i$, $\thickmuskip=2mu \medmuskip=2mu \bm{R}_i\in\mathbb{R}^{d/r\times d/r}$ are both orthogonal. When the block size $r$ is set to $1$, the block-diagonal orthogonal matrix simplifies to a fully unconstrained orthogonal matrix. The parameter count for an orthogonal matrix of size $\thickmuskip=2mu \medmuskip=2mu d\times d$ is given by $\thickmuskip=2mu \medmuskip=2mu d(d-1)/2$, resulting in an overall computational complexity of $\mathcal{O}(d^2)$. In contrast, an $r$-block diagonal orthogonal matrix comprises $\thickmuskip=2mu \medmuskip=2mu d(d/r-1)/2$ parameters, leading to complexity $\thickmuskip=2mu \medmuskip=2mu \mathcal{O}(d^2/r)$. By tying, or sharing, the block matrices across different blocks (i.e., imposing $\thickmuskip=2mu \medmuskip=2mu \bm{R}_i=\bm{R}_j,\forall i\neq j$), the parameter count can be reduced further, resulting in complexity $\mathcal{O}(d^2/r^2)$. Notably, each strategy for reducing the parameter count preserves the orthogonality of the resulting $\bm{R}$ matrix, thereby maintaining the same hyperspherical energy.

Comparing parameter efficiency of OFT and LoRA, for LoRA with rank parameter $r'$, the number of learnable parameters is given by $\thickmuskip=2mu \medmuskip=2mu r'(d+n)$. If both $r$ and $r'$ scale linearly with $d$ (for instance, $\thickmuskip=2mu \medmuskip=2mu r=r'=\alpha d$ with constant $0<\alpha\leq 1$), LoRA’s parameter complexity becomes $\mathcal{O}(d^2+dn)$, whereas for OFT, it is reduced to $\mathcal{O}(d)$. To make this more concrete, consider a weight matrix of dimensions $\thickmuskip=2mu \medmuskip=2mu 128\times 128$: for $r'=8$, LoRA requires $2,048$ trainable parameters, whereas for $r=8$ (without block sharing), OFT needs only $960$ trainable parameters.

\subsection{Constrained Orthogonal Finetuning}

It is possible to further restrict the expressiveness of standard OFT by forcing the finetuned parameters to remain within a limited vicinity of the original pretrained model. Constrained Orthogonal Finetuning (COFT) achieves this by enforcing the following during the forward computation:

\begin{equation}\label{eq:coft_general}
    \footnotesize
    \bm{z}=\bm{W}^\top\bm{x}=(\bm{R}\cdot\bm{W}^0)^\top\bm{x},~~~\text{s.t.}~\bm{R}^\top\bm{R}=\bm{R}\bm{R}^\top=\bm{I},~\norm{\bm{R}-\bm{I}}\leq \epsilon
\end{equation}

This introduces both an orthogonality constraint and an $\epsilon$-proximity constraint to the identity matrix. Orthogonality can be realized through Cayley parameterization, as detailed in Section~\ref{sect:eff_ortho}. Nevertheless, directly encoding the $\epsilon$-deviation constraint within the Cayley parameterized form is a complex task. To better understand the structure, one can approximate the Cayley transform $\thickmuskip=2mu \medmuskip=2mu \bm{R}=(\bm{I}+\bm{Q})(I-\bm{Q})^{-1}$ using the Neumann series, which gives $\thickmuskip=2mu \medmuskip=2mu \bm{R}\approx\bm{I}+2\bm{Q}+\mathcal{O}(\bm{Q}^2)$ under suitable assumptions on $\bm

series is convergent with respect to the operator norm). Consequently, the condition $\thickmuskip=2mu \medmuskip=2mu\|\bm{R}-\bm{I}\|\leq\epsilon$ can be incorporated within the Cayley transform, leading to an equivalent restriction of the form $\thickmuskip=2mu \medmuskip=2mu \|\bm{Q}-\bm{0}\|\leq\epsilon'$, where $\bm{0}$ stands for the zero matrix, and $\epsilon'$ serves as a new hyperparameter for tolerable error (generally distinct from $\epsilon$). This alternative constraint imposed on $\bm{Q}$ can be readily satisfied using projected gradient descent techniques. To initialize the orthogonal matrix $\bm{R}$ to the identity, we simply set $\bm{Q}$ to be the zero matrix at the start. COFT unifies two explicit regularization objectives: limiting the change in hyperspherical energy and enforcing proximity to the original pretrained weights. While the latter is frequently used in standard finetuning as an implicit constraint, COFT formalizes it as an explicit regularization. Despite OFT’s strong empirical performance, our findings indicate that COFT further enhances stability during finetuning due to this explicit deviation penalty. Figure~\ref{fig:eps_coft} illustrates the influence of $\epsilon$ on COFT’s effectiveness. It is evident that $\epsilon$ modulates the degree of adaptation: higher values allow COFT-finetuned models to approach the behavior of OFT-finetuned counterparts, whereas lower values constrain the COFT solution to closely mimic the pretrained text-to-image diffusion model.

\subsection{Re-scaled Orthogonal Finetuning}

We introduce a straightforward enhancement to OFT by allowing the model to learn a scaling factor for each neuron in addition to the orthogonal transformation. The rationale behind this modification is that scaling the activations of individual neurons leaves their hyperspherical energy unaffected (as the normalization removes the magnitude component). Specifically, we formulate the forward computation as $\thickmuskip=2mu \medmuskip=2mu\bm{z}=(\bm{R}\bm{W}^0\bm{D})^\top\bm{x}$\footnote{\textbf{Errata}: The camera-ready NeurIPS version incorrectly specified the forward calculation for re-scaled OFT as $\thickmuskip=2mu \medmuskip=2mu\bm{z}=(\bm{D}\bm{R}\bm{W}^0)^\top\bm{x}$; however, the original code was correct and results in Appendix~\ref{app:rescaled} remain unaffected.} where $\thickmuskip=2mu \medmuskip=2mu \bm{D}=\text{diag}(s_1,\cdots,s_n)\in\mathbb{R}^{n\times n}$ denotes a learnable diagonal matrix whose diagonal entries $s_1,\ldots,s_n$ are constrained to be positive. While standard OFT (see Eq.~\eqref{eq:oft_general}) only involves training $\bm{R}$, the re-scaled variant jointly optimizes both the diagonal scaling matrix $\bm{D}$ and the orthogonal $\bm{R}$. This adjustment grants OFT enhanced modeling capacity with only a marginal increase in parameter count. In our experiments, we focus on the canonical OFT framework to better demonstrate the impact of orthogonal transforms, though our investigations indicate that the re-scaled variant consistently delivers superior performance (refer to Appendix~\ref{app:rescaled}).

\section{Intriguing Insights and Discussions}

\textbf{OFT is architecture-agnostic}. The OFT methodology is, in principle, adaptable to any neural network architecture. Within Transformer networks, LoRA is most frequently applied to attention weight matrices~\cite{hulora2022}. For equitable comparison with LoRA, our experiments restrict OFT to only update attention parameters. OFT is not limited to dense layers; it is especially natural for convolutional layers, as its block-diagonal structure admits unique interpretations in this context—unlike LoRA. When the number of blocks matches the input channel count, each block operates on a separate neuron channel, analogous to depth-wise convolutions~\cite{chollet2017xception}. Conversely, if all blocks are identical and shared, the orthogonal transformation is applied uniformly across all channels. A preliminary exploration of applying OFT to convolutional finetuning is included in Appendix~\ref{app:convolution}.

\textbf{Connection to LoRA}. LoRA introduces a low-rank update to the pretrained weight matrix, which helps to avoid significant alterations to the original weights. In contrast, OFT regulates the transformation applied to the pretrained weight matrix by constraining it to be orthogonal (thus, full-rank), thereby ensuring the integrity of the information learned during pretraining is maintained. The forward computation for OFT can be reformulated as $\thickmuskip=2mu \medmuskip=2mu \bm{z}=(\bm{R}\bm{W}^0)^\top\bm{x}=(\bm{W}^0+(\bm{R}-\bm{I})\bm{W}^0)^\top\bm{x}$, in which $\thickmuskip=2mu \medmuskip=2mu (\bm{R}-\bm{I})\bm{W}^0$ parallels LoRA’s low-rank modification. Given that $\bm{W}^0$ is usually full-rank, when $\thickmuskip=2mu \medmuskip=2mu \bm{R}-\bm{I}$ has low rank, OFT effectively carries out a low-rank adjustment akin to LoRA. OFT leverages a diagonal block size $r$, akin to LoRA’s rank parameter $r'$, to limit the number of trainable variables. Notably, LoRA and OFT embody two fundamentally different mechanisms for parameter efficiency: LoRA capitalizes on a low-rank parameterization, whereas OFT leverages block-diagonal orthogonal transformations to induce sparsity in the learnable parameters.

\textbf{Why OFT converges faster}. Firstly, as illustrated in Figure~\ref{fig:toy_ae}, adjusting neuron orientations most effectively alters semantic content, aligning directly with the objective of OFT. Furthermore, OFT’s training process can be understood as adapting neurons while constraining them to remain on a smooth hypersphere, leading to a more tractable and favorable optimization geometry. Experimental results supporting this have been demonstrated in \cite{liu2021orthogonal}.

\textbf{Why not minimize hyperspherical energy}. A salient distinction from \cite{liu2021orthogonal} lies in the fact that our approach does not seek to minimize hyperspherical energy. For classification tasks, it is advantageous for neurons to be non-redundant; however, the minimization of hyperspherical energy produces a configuration where neurons are distributed uniformly across the hypersphere. This is not a suitable goal for finetuning since such uniformity may disrupt or erase valuable pretraining knowledge.

\textbf{Balancing flexibility and regularization in finetuning}. Our findings reveal a fundamental tension between model flexibility and regularization during finetuning. Conventional finetuning allows for maximal adaptation but at the cost of stability, frequently resulting in model degradation. In contrast, OFT—while conceptually straightforward—achieves a compromise by maintaining the angular relationships between neurons, thereby mediating between adaptability and imposed structure. Moreover, the block-diagonal parametrization serves as a more stringent regularizer for the orthogonal transformation matrix.

\textbf{No extra inference cost}. In contrast to approaches like ControlNet, OFT incurs zero additional inference latency. After training, we simply compose the learned orthogonal matrix $\bm{R}$ with the original parameter matrix $\bm{W}^0$, yielding $\thickmuskip=2mu \medmuskip=2mu \bm{W}=\bm{R}\bm{W}^0$, a transformation indistinguishable in computational cost from the pretrained model. Hence, inference throughput remains unchanged.

\section{Experiments and Results}

\textbf{Experimental protocols}. Our empirical evaluation employs Stable Diffusion v1.5~\cite{rombach2022high} as the foundational text-to-image architecture. For consistent comparison, we randomly select generated samples from all evaluated approaches. In subject-driven synthesis, our setup is aligned with the procedures outlined in DreamBooth~\cite{ruiz2023dreambooth}. For experiments on controllability, protocols from ControlNet~\cite{zhang2023adding} and T2I-Adapter~\cite{mou2023t2i} are followed. To ensure comparability with LoRA, we restrict the application of OFT or COFT to the identical network layers as those used by LoRA. Further details and expanded results are located in Appendix~\ref{app:settings}.

\subsection{Subject-driven Generation}

\textbf{Configuration}. Baseline comparisons are conducted against DreamBooth~\cite{ruiz2023dreambooth} and LoRA~\cite{hulora2022}. All evaluated models are trained with the same objective function prescribed by DreamBooth. Hyperparameter choices for DreamBooth and LoRA are drawn from their respective source publications, using optimal configurations as reported. Expanded experimental data can be found in Appendix~\ref{app:settings},\ref{app:coft_oft},\ref{app:more_qualitative},\ref{app:failure}.

\textbf{Finetuning stability and convergence}. Our initial investigation focuses on assessing the stability during finetuning and the rate of convergence for DreamBooth, LoRA, OFT, and COFT. The findings, visualized in Figure~\ref{fig:teaser} and Figure~\ref{fig:db_convergence}, indicate that both COFT and OFT enable the diffusion model to be finetuned in a robust manner. By the 400-iteration mark, both DreamBooth and the OFT-based approaches deliver effective subject control, whereas LoRA demonstrates deficiencies in maintaining subject identity. At 2000 iterations, DreamBooth produces degraded outputs (collapsed images), and LoRA incorrectly synthesizes a yellow fur in place of a yellow shirt. Conversely, OFT and COFT consistently maintain stable generation quality and subject control across iterations. These observations highlight the rapid yet reliable convergence characteristics of the OFT framework when used for subject-driven synthesis. Notably, the low sensitivity of OFT and COFT to the precise number of finetuning iterations is advantageous, as it simplifies the process by removing the need for fine-tuning the iteration count per subject. Both OFT and COFT support the use of a relatively high number of iterations without careful adjustment. When using COFT with an appropriate choice of $\epsilon$, selecting the learning rate and iteration count becomes a trivial task.

\textbf{Quantitative comparison}. Consistent with \cite{ruiz2023dreambooth}, we carry out a quantitative analysis assessing subject fidelity (via DINO~\cite{caron2021emerging} and CLIP-I~\cite{radford2021learning}), alignment with textual prompts (CLIP-T~\cite{radford2021learning}), and output diversity (LPIPS~\cite{zhang2018unreasonable}). Specifically, CLIP-I is determined by computing the mean pairwise cosine similarity between CLIP embeddings of synthesized and reference images. The DINO metric follows an analogous protocol, but leverages ViT S/16 DINO representations. CLIP-T quantifies the mean cosine similarity between the CLIP embeddings of text prompts and corresponding generated images. For diversity, we report the average LPIPS cosine similarity for images generated using identical subjects and text prompts. According to Table~\ref{table:dreambooth_final}, both COFT and OFT substantially outperform DreamBooth and LoRA on DINO and CLIP-I, while also achieving either similar or slightly superior results for prompt fidelity and diversity. Each experimental protocol involves 30 repeated finetuning runs per method with different random seeds, reducing the effect of stochastic variation. Overall, the results affirm that OFT not only improves convergence stability, but also achieves enhanced final generation performance in a consistent manner.

\textbf{Qualitative comparison}. To facilitate a more intuitive grasp of the advantages provided by OFT, we present several randomly selected samples for subject-driven generation in Figure~\ref{fig:dreambooth_final}. In order to ensure fairness in comparison, each example is produced by the same finetuned model across all methods, and no separate prompt tuning is conducted for any final image. For each approach, the version of the model achieving the highest validation CLIP scores is chosen. Upon inspecting the outcomes illustrated in Figure~\ref{fig:dreambooth_final}, it is clear that OFT and COFT are both highly effective at retaining the semantic identity of the subject, whereas LoRA frequently struggles with subject fidelity (\emph{e.g.}, LoRA entirely misrepresents the subject in the bowl instance). Furthermore, both OFT and COFT exhibit superior responsiveness to textual prompts, in contrast to DreamBooth, which, while generally successful at subject preservation, often fails to comply with the specified prompts (\emph{e.g.}, the first row of the bowl case). These qualitative results underscore that the OFT framework excels in maintaining both controllability and subject fidelity concurrently. Additionally, OFT demonstrates robustness to the number of iterations—performing reliably even as iteration counts increase—unlike DreamBooth and LoRA. Further qualitative samples are provided in Appendix~\ref{app:more_qualitative}. A human evaluation reported in Appendix~\ref{app:human} also corroborates the superior performance of OFT.

\subsection{Controllable Generation}

\textbf{Settings}. For baseline comparisons, we employ ControlNet~\cite{zhang2023adding}, T2I-Adapter~\cite{mou2023t2i}, and LoRA~\cite{hulora2022}. The main text examines three demanding controllable generation tasks: translating Canny edge maps to images~(C2I) using the COCO dataset~\cite{lin2014microsoft}, converting segmentation maps to images~(S2I) with the ADE20K dataset~\cite{zhou2017scene}, and generating facial images from landmarks~(L2F) on the CelebA-HQ dataset~\cite{karras2018progressive,xia2021tedigan}. All methods are utilized to finetune Stable Diffusion~(SD) v1.5 across these datasets for 20 epochs. Additional results can be found in Appendix~\ref{app:more_qualitative},\ref{app:more_control_task},\ref{app:failure}.

\textbf{Convergence}. To assess convergence rates, we examine the performance of ControlNet, T2I-Adapter, LoRA, and COFT on the L2F benchmark through both quantitative metrics and qualitative observations. For quantitative assessment, we calculate the mean $\ell_2$ distance separating the predicted face landmarks from the control landmarks. Figure~\ref{fig:face_convergence} presents how face landmark error varies as the models are finetuned over an increasing number of epochs. Notably, both COFT and OFT display markedly accelerated convergence compared to the baselines; for instance, LoRA requires 20 epochs to reach the error attained by OFT at merely the 8-th epoch. It is important to highlight that OFT and COFT retain a trainable parameter count on par with LoRA (and substantially below that of ControlNet), yet achieve much swifter convergence. The accelerated learning capacity of OFT is further supported by findings in Figure~\ref{fig:teaser}, particularly the rightmost example, which demonstrates that OFT is notably more data-efficient than both ControlNet and LoRA, converging robustly with just 5\% of the ADE20K dataset. For the qualitative analysis, we primarily compare OFT, COFT, and ControlNet, as ControlNet attains a comparable landmark error. In Figure~\ref{fig:face_convergence_qual}, we see that both OFT and COFT undergo smooth and reliable convergence, with generated facial poses incrementally aligning with the control landmarks over epochs. In contrast, ControlNet exhibits abrupt emergence of controllability after the 8-th epoch, corroborating the sudden convergence pattern reported in \cite{zhang2023adding}. The consistent and gradual convergence observed in our OFT framework facilitates its practical adoption, as it results in smoother and more foreseeable training dynamics, thereby simplifying the process of hyperparameter selection.

\textbf{Quantitative comparison}. To assess controllable generation, we propose a control consistency metric, which operates by extracting the control signal from the synthesized image and contrasting it with the initial input control. Specifically, for the C2I task, Intersection-over-Union (IoU) and F1 scores are reported. For the S2I task, we evaluate mean IoU as well as mean and overall accuracy, while for the L2F task, the average $\ell_2$ distance between the target control landmarks and those predicted is calculated. Comprehensive descriptions of these metrics can be found in Appendix~\ref{app:settings}. All compared methods utilize optimally tuned hyperparameters. As shown in Table~\ref{table:control_gen}, OFT and COFT consistently deliver markedly superior and more precise control in comparison to competing approaches. Adapter-based methods, such as T2I-Adapter and ControlNet, not only converge more slowly but also underperform in terms of final control quality. Notably, LoRA achieves higher results than ControlNet on the S2I task but lags behind on C2I and L2F tasks. Generally, we observe a limited performance upper bound for LoRA, even with extensive hyperparameter tuning. Conversely, the OFT framework’s capability has not yet plateaued; empirically, its performance improves with an increased number of trainable parameters. Importantly, our framework for quantitative evaluation in the context of controllable generation represents a key novel contribution, offering an accurate means to gauge control fidelity in finetuned models, subject to the accuracy constraints of the available segmentation or detection backbone.

\textbf{Qualitative comparison}. We further conduct a qualitative evaluation of OFT and COFT in comparison to prevailing state-of-the-art approaches, namely ControlNet, T2I-Adapter, and LoRA. As illustrated by the randomly generated results in Figure~\ref{fig:control_final}, both OFT and COFT produce images characterized by high realism and fidelity, while simultaneously maintaining precise control over the output. Within the S2I task, it becomes apparent that LoRA is unable to synthesize images that correspond to the provided segmentation maps, whereas ControlNet, OFT, and COFT successfully enforce such spatial constraints. Unlike ControlNet, both OFT and COFT deliver images with enhanced detail, realism, and plausible geometrical configurations, all while utilizing significantly fewer model parameters. Regarding the C2I task, OFT and COFT consistently generate semantically coherent images even from coarse Canny edge inputs, outperforming both T2I-Adapter and LoRA in terms of controllability and quality. In the L2F task, our proposed methods exhibit superior accuracy in pose guidance, especially in scenarios involving challenging facial orientations. Across these three controlled generation benchmarks, the qualitative results suggest that OFT and COFT surpass the baseline models, underscoring the robust controllable generation capability of our OFT framework. For an expanded set of qualitative results spanning all three control settings, refer to Appendix~\ref{app:control_exp}; additionally, Appendix~\ref{app:more_control_task} demonstrates the versatility of OFT across a broader array of control tasks, including dense pose to human body, sketch to image, and depth to image transformations.

\section{Concluding Remarks and Open Problems}

Inspired by the crucial role of angular relationships among neurons in defining visual semantics, we have introduced a straightforward yet robust approach—orthogonal finetuning—for the task of controlling text-to-image diffusion models. Our investigation focuses on two key text-to-image scenarios: subject-driven generation and controllable generation. When compared with prior approaches, OFT yields enhanced controllability and greater finetuning stability while utilizing a reduced set of tunable parameters. Notably, OFT does not incur any extra computational cost during inference, supporting its suitability for practical deployment.

Several intriguing open questions emerge from OFT. Firstly, maintaining orthogonality via Cayley parametrization necessitates a matrix inverse, which can hinder the scalability of OFT. Although this issue is partly addressed through the use of block diagonal parametrization, developing faster differentiable solutions for the matrix inversion step remains unresolved. Secondly, OFT exhibits promising compositional capabilities, as the orthogonal matrices generated by separate OFT finetuning procedures can be multiplied, resulting in a product that remains orthogonal. Investigating whether these orthogonal matrix products successfully retain the learned information from all involved downstream tasks presents an important avenue for future research. Lastly, the model's parameter efficiency is primarily influenced by its block diagonal configuration, which inevitably brings about certain biases and constraints on adaptability. Developing strategies to further enhance parameter efficiency while minimizing bias stands as a key open challenge.